% GAT-26 — BraTS 2026 Task 3 (BraTS-GoAT) short paper.
% Springer LNCS format, 8–10 pages (excluding references), per the live-verified submission rules
% (Synapse wiki/639582, re-verified 2026-07-25T12:52Z).
%
% BUILD: this uses the Springer LNCS class `llncs.cls`, which is copyrighted by Springer and is NOT
% committed here. Download the official LNCS template (https://www.springer.com/gp/computer-science/lncs)
% into this directory, then compile with the command in paper/README.md.
%
% CONTENT POLICY: only VERIFIED committed evidence is populated. Everything not yet measured/decided
% is a conspicuous placeholder marked \ownerinput{...} (owner input required) or \pending{...}
% (pending A10G-2 / official validation / hidden test). NEVER fabricate or extrapolate final results.
%
% REQUIRED SECTIONS (wiki/639582): Abstract (no citations), keywords, Introduction, Methods, Results,
% Discussion, Acknowledgements (optional), References (must include the Challenge-Rules citations).

\documentclass[runningheads]{llncs}
\usepackage{graphicx}
\usepackage{booktabs}
\usepackage{amsmath}
\usepackage[hidelinks]{hyperref}
\usepackage{xurl}   % allow long URLs to break across lines (the LNCS text block is narrow)
\usepackage{xcolor}

% Conspicuous fail-closed placeholders (must be resolved before submission).
\newcommand{\ownerinput}[1]{\textbf{\textcolor{red}{[OWNER\_INPUT\_REQUIRED: #1]}}}
\newcommand{\pending}[1]{\textbf{\textcolor{red}{[PENDING: #1]}}}

\begin{document}

% Working title. The owner supplies the final title; the descriptive subtitle below is the
% system/task description and may be kept, shortened, or replaced.
\title{\ownerinput{final title} --- GAT-26: A Generalization-Aware Tri-region 3D nnU-Net
Ensemble for Cross-Tumor Brain Tumor Segmentation}
% `runningheads' requires a short running title; without it LNCS emits
% "Title too long for running head".
\titlerunning{GAT-26: Tri-region 3D nnU-Net Ensemble for BraTS-GoAT 2026}

% The challenge runs a light peer review in which participants review 2–3 papers (wiki/639582,
% re-verified 2026-07-25). No anonymization requirement was found on any challenge page, so the
% manuscript is prepared NON-ANONYMOUSLY (authors named). The exact OpenReview venue setting remains
% an OWNER-VERIFICATION item before submission.
\author{\ownerinput{Author One} \and \ownerinput{Author Two}}
\authorrunning{\ownerinput{Short author list}}
\institute{\ownerinput{Affiliation(s) and contact emails}}

\maketitle

\begin{abstract}
% Abstract: no citations (per wiki/639582 section rules).
We describe GAT-26, our submission to the BraTS 2026 Generalizability Across Tumors (BraTS-GoAT)
subchallenge (Task 3): pre-operative, cross-tumor 3D brain tumor segmentation of the enhancing tumor
(ET), tumor core (TC), and whole tumor (WT) regions. GAT-26 is built on the official nnU-Net v2.8.1
Residual-Encoder U-Net with region-based training, trained only on GoAT-provided data from a recorded
random initialization with no external or pretrained weights. Architecture selection between the
ResEnc-M and ResEnc-L presets was made by a pre-registered, result-blind policy on a leakage-safe
fold-0 screen; the policy selected ResEnc-M. We then trained the complete five-fold cross-validation
and reconciled a strictly out-of-fold (OOF) prediction set covering every one of the 1,351 training
cases exactly once, scored in a single pass by the official evaluator with zero errors. Two further
pre-registered, calibration-then-confirmation audits of the inference policy — one over checkpoint
choice, mirroring test-time augmentation and enhancing-tumor component cleanup, and one over
checkpoint-weight averaging (``model soups'') — both retained the frozen baseline policy, which we
report as evidence that the released configuration is not the product of post-hoc tuning. The final
system is a sequential five-checkpoint mean-probability ensemble with a fixed 0.5 threshold and a
hierarchy-safe reconstruction, executed under enforced-determinism settings inside a zero-network
clean-room container. The five trained models and the frozen inference policy exist and are fixed;
the final container image, its genuine-A10G runtime and memory acceptance, official validation
performance, and hidden-test performance are not yet available and are reported as pending.
\keywords{Brain tumor segmentation \and Generalization \and nnU-Net \and Residual encoder \and
Cross-validation \and BraTS.}
\end{abstract}

\section{Introduction}
The Brain Tumor Segmentation (BraTS) benchmark has, over more than a decade, standardized the
evaluation of automated glioma segmentation from multi-parametric
MRI~\cite{menze2015brats,bakas2017advancing,baid2021brats}, and has since expanded into a family of
cohort-specific subchallenges. These span intracranial meningioma~\cite{bratsmen2023},
brain metastases~\cite{bratsmet2023}, pediatric tumors~\cite{bratsped2023} and an
under-represented Sub-Saharan Africa patient population~\cite{bratsssa2023}, and are
scored by containerized evaluation infrastructure in the spirit of federated
benchmarking~\cite{karargyris2023medperf}.

Cross-tumor generalization is the central difficulty of the BraTS-GoAT task: a single model must
segment ET/TC/WT across these heterogeneous adult and pediatric tumor types and acquisition sources
without per-cohort tuning. GAT-26 targets this with a deliberately conservative, reproducible pipeline: a
strong region-based nnU-Net backbone, a leakage-safe group-level cross-validation design,
pre-registered decision rules that cannot be edited after seeing results, and a release-frontier
container constrained to the official Task-3 runtime budget.

Our emphasis is methodological discipline rather than architectural novelty. Every decision that
could be tuned — architecture preset, checkpoint choice, test-time augmentation, post-processing,
and checkpoint averaging — was placed behind a frozen policy with a held-out confirmation stage, and
in every case the policy retained the simpler baseline. We therefore report a cross-validated
characterization of the system together with an explicit statement of what has \emph{not} been
measured, rather than a competitiveness claim.

\section{Methods}

\subsection{Data and Task}
We use only data supplied through the 2026 BraTS-GoAT subchallenge website, which draws on the
BraTS data lineage and its expert-annotated public
collections~\cite{menze2015brats,bakas2017advancing,bakas2017gbm,bakas2017lgg,baid2021brats}. The training release
comprises \textbf{1{,}351} labeled cases (verified in our data audit: all complete and geometry-valid,
integer labels only, zero label anomalies). Labels follow the BraTS convention and the evaluated
regions are the standard nested triple:
\begin{equation*}
\begin{gathered}
\{0{=}\text{background},\;1{=}\text{NCR},\;2{=}\text{ED},\;3{=}\text{ET}\},\\[2pt]
\mathrm{ET}=[3],\qquad \mathrm{TC}=[1,3],\qquad \mathrm{WT}=[1,2,3].
\end{gathered}
\end{equation*}
Validation and test ground-truth
labels are withheld by the organizers; we therefore do not train on validation data and never use it
for tuning. Per-case cohort labels (e.g. adult glioma, meningioma, metastasis, pediatric, Sub-Saharan
Africa) were \emph{not} reliably derivable from the available file naming or metadata, and we did not
guess them; consequently no per-cohort sizes or cohort identities are claimed, and cohort-stratified
analysis is unavailable in this work.

\subsection{Leakage-safe Grouping and Five-fold Split}
To avoid optimistic leakage, we build a deterministic \emph{group-level} five-fold split (seed
$21072026$). A duplicate/near-duplicate analysis found no confirmed same-subject pairs; the 1{,}351
cases are treated as singleton groups, and we explicitly acknowledge that residual same-subject risk
cannot be proven to be exactly zero. The split is reproducible: two independent runs produced
byte-identical assignments. Fold sizes are $[271,270,270,270,270]$; every fold contains ET-absent and
TC-absent cases, exercising the empty-region edge cases the evaluator must handle. The split was
frozen before any training and was never modified.

\subsection{Preprocessing}
Preprocessing is the standard nnU-Net~\cite{isensee2021nnunet} v2.8.1 pipeline for the four input modalities (T1n, T1c, T2w,
T2f) in a fixed channel order, using the dataset fingerprint and plans generated by nnU-Net planning.
No external normalization statistics or external data are introduced.

\subsection{ResEnc-M Architecture and Training}
The backbone is the official nnU-Net~\cite{isensee2021nnunet} v2.8.1 \texttt{ResidualEncoderUNet}
(ResEnc) preset~\cite{isensee2024nnunetrevisited} in the \texttt{3d\_fullres} configuration with
region-based training (the network predicts the ET/TC/WT regions directly). \textbf{All trainable parameters begin from a recorded random initialization; no
pretrained, foundation, or self-supervised weights are used}, and unexpected weight files are treated
as a release-blocking error. Each fold is trained for 1000 epochs with the native region-based loss,
from an independent random initialization, with no cross-fold resume and no warm starting.
\textbf{All five folds (0--4) are complete}; fold~0 is the model produced by the architecture-selection
screen and was never retrained.

\subsection{Pre-registered Architecture Selection}
Architecture selection between ResEnc-M and ResEnc-L was governed by a \textbf{pre-registered,
result-blind policy fixed before any fold result was seen} (and unchanged thereafter). The internal
award utility is the mean fractional rank over six equally-weighted components (ET/TC/WT $\times$
DSC/HD95); DSC is higher-better and HD95 lower-better; ties resolve to the smaller, cheaper model
(ResEnc-M). A paired subject-level bootstrap (seed $21072026$, $10{,}000$ resamples) and a set of
pre-declared noninferiority margins guard the decision. ResEnc-L is preferred only if it shows a
\emph{meaningful} fold-0 benefit (award-rank gain $\ge 1/6$, a bootstrap interval favoring L, and all
noninferiority gates passing); otherwise the policy selects ResEnc-M. This screen selects the
architecture only; it is \emph{not} a competitiveness claim. Section~\ref{sec:metric} records an
important caveat about the relationship between this internal utility and the official ranking
metrics.

\subsection{Inference Ensemble and Hierarchy-safe Reconstruction}
Final inference uses a \textbf{sequential} five-checkpoint mean-probability ensemble: one fold model is
resident at a time and freed before the next, so peak memory is that of a single model. Region
probabilities are thresholded at a fixed $0.5$ (no learned threshold, no learned ensemble weight, no
test-time augmentation, no connected-component filtering, no presence gate) and reconstructed with a
hierarchy-safe projection enforcing $\mathrm{WT}\supseteq\mathrm{TC}\supseteq\mathrm{ET}$. Outputs
restore the exact source geometry (shape, spacing, origin, direction) and are integer-labeled
$\{0,1,2,3\}$. Inference runs under enforced-determinism settings (fixed seeds, deterministic
cuDNN/cuBLAS kernels). We measured reproducibility rather than assuming it: across two independent
24-case runs, $18$ of $2.14\times10^{8}$ output voxels differed ($8.4\times10^{-8}$, affecting 9 of
24 cases by 1--5 voxels each), while a repeated 4-case subset was bit-identical across three runs.
The runner is therefore \emph{near}-deterministic but not bit-exact in general: PyTorch's
deterministic-algorithm enforcement is applied in \texttt{warn\_only} mode, so the few operations
lacking a deterministic CUDA kernel may still flip borderline voxels across the $0.5$ threshold. The
magnitude is far below any effect on the reported metrics, and every output remains label-valid,
hierarchy-consistent and geometry-exact.

\subsection{Frozen Inference-policy Audits}
\label{sec:audits}
Two bounded audits tested whether the frozen inference policy could be improved, each using a
\emph{calibration-then-confirmation} design on disjoint folds, the frozen official evaluator, the same
six-component utility, and the same paired bootstrap, with the baseline winning every tie:

\begin{itemize}
  \item \textbf{Audit~A (checkpoint / TTA / post-processing).} Candidates: the best-validation
        checkpoint; mirroring test-time augmentation over the axes recorded in the
        checkpoint metadata; both combined; and enhancing-tumor component cleanup at three
        pre-declared physical-volume thresholds. Calibration used folds 0--2 ($n=811$). \textbf{No candidate produced a robust
        improvement}; test-time augmentation was point-positive but its bootstrap interval did not
        exclude a tie, and component cleanup traded ET Dice against ET boundary distance for an exact
        six-component tie. The baseline was retained and the locked confirmation folds were never
        opened.
  \item \textbf{Audit B (checkpoint soups).} Two pre-declared weight-averaging candidates per fold,
        $S_1=0.75\,\theta_{\mathrm{final}}+0.25\,\theta_{\mathrm{best}}$ and
        $S_2=0.50\,\theta_{\mathrm{final}}+0.50\,\theta_{\mathrm{best}}$, constructed with strict
        structural guards (identical key sets and shapes, floating-point tensors only, non-floating
        buffers required to be bitwise identical, no optimizer or scheduler state carried over) and
        loaded strictly. Calibration on folds 0--2 ($n=811$). \textbf{Both candidates were worse than
        the baseline on all six award components} (mean-fractional-rank gain $-1.000$ each), the paired
        bootstrap favoured the baseline, and the 0.75/0.25 soup additionally failed the pre-declared
        95th-percentile boundary-distance noninferiority gate. The baseline was retained and the
        locked confirmation folds were never opened.
\end{itemize}

The purpose of reporting these audits is not the candidates themselves but the \emph{negative
result}: the released configuration is the pre-registered baseline, not a post-hoc selection.

\subsection{Containerization and Reproducibility}
The submission target is a Linux/AMD64, zero-network clean-room container designed for the official
Task-3 budget (NVIDIA A10G, 24\,GB VRAM, 16 vCPU, \texttt{--memory} 48\,GiB, \texttt{--shm-size}
16\,GiB, CUDA $\le 13.0$, 12\,h total). Dependencies are fully pinned and baked at build time,
\texttt{/input} is mounted read-only, and \texttt{/output} is a flat directory with exactly one
\texttt{.nii.gz} per case whose name ends in the five-digit case identifier. Strict modality discovery
and fail-closed input validation reject missing, duplicate, unknown, or unreadable modalities, invalid
folder names, and output-name collisions \emph{before any output is written}. The runner refuses to
start unless exactly five distinct fold checkpoints are present.

% Public source-code URL, filled in Stage G79-P after the sanitized export was published and
% independently verified by unauthenticated clone. Only this field was filled; every other
% \ownerinput{} / \pending{} placeholder remains open.
Code availability: the source code for this system is publicly available under the Apache License
2.0 at \url{https://github.com/Neethan-hub/gat26-brats-2026}. The repository contains the inference
runner and its container definition, the split/selection/training and evaluation tooling,
configuration, tests, and this paper's source. It contains no images, labels, model checkpoints, or
predictions; the challenge data must be obtained independently from the organizers under their own
access terms.

\section{Results}

\subsection{Architecture-selection Screen (fold 0)}
\textbf{This subsection reports the fold-0 architecture-selection screen only.} On the frozen fold-0
validation set (271 cases), the official evaluator (BraTS-evaluation 0.0.8, GoAT configuration)
produced the region means in Table~\ref{tab:fold0}. Legitimate zero-true-positive $+\infty$ HD95
values are converted to the organizer penalty of 373\,mm and retained in every mean and denominator
(all six components use the full 271 cases).

\begin{table}[t]
\centering
\caption{Fold-0 selection screen (271 validation cases; means). Architecture-selection evidence, not
a competitiveness result. DSC higher-better; HD95 (mm) lower-better.}
\label{tab:fold0}
\begin{tabular}{lccc}
\toprule
Model & ET (DSC / HD95) & TC (DSC / HD95) & WT (DSC / HD95) \\
\midrule
ResEnc-M (selected) & 0.859 / 14.17 & \textbf{0.914 / 5.96} & \textbf{0.934 / 3.89} \\
ResEnc-L            & \textbf{0.861 / 11.19} & 0.912 / 6.08 & 0.932 / 4.11 \\
\bottomrule
\end{tabular}
\end{table}

Under the frozen policy, ResEnc-M wins four of the six components; the paired bootstrap yields a rank
difference that does not favor ResEnc-L (95\% interval $[-0.667,\,+0.667]$), and ResEnc-L fails the
pre-declared 5th-percentile DSC and 95th-percentile HD95 tail-noninferiority gates (DSC 5th pct
$0.649$ vs $0.637$; HD95 95th pct $15.96$ vs $18.06$\,mm). The policy therefore \textbf{selects
ResEnc-M}.

\subsection{Five-fold Cross-validated Out-of-fold Performance}
All five ResEnc-M folds were trained and audited to completion. We assembled a strictly out-of-fold
prediction set: each case is predicted exactly once, by the fold model that did \emph{not} train on
it. The union covers \textbf{1{,}351 cases, each exactly once}, with no missing case, no duplicate,
and no fold overlap; each fold's realized predictions equal its frozen validation membership. The
pooled set was scored in a single pass by the official evaluator: \textbf{$n=1{,}351$, zero errors},
with all six component denominators equal to 1{,}351 (Table~\ref{tab:oof}).

\begin{table}[t]
\centering
\caption{Five-fold cross-validated \textbf{out-of-fold} performance (1{,}351 cases, each predicted
once by a model that did not train on it; official evaluator, $n=1{,}351$, 0 errors).
\textbf{This is cross-validation evidence, not validation-leaderboard or hidden-test performance.}}
\label{tab:oof}
\begin{tabular}{lccc}
\toprule
Region & DSC & HD95 (mm) & denominator \\
\midrule
ET & 0.865 & 13.39 & 1351 \\
TC & 0.915 & 6.75  & 1351 \\
WT & 0.927 & 6.34  & 1351 \\
\bottomrule
\end{tabular}
\end{table}

Tail behaviour, which matters more than the mean for a generalization task, is characterized by a
5th-percentile DSC of $0.666$ and a 95th-percentile HD95 of $17.72$\,mm. Across the
$1351\times3=4053$ region instances, 3{,}987 scored normally, 49 were zero-true-positive penalties,
8 were disjoint penalties, and 9 were true negatives where the region was legitimately absent from
both reference and prediction. One fold's evaluation initially failed on a tooling defect (an
evaluator denominator was not passed explicitly); the defect was fixed, the preserved predictions
were re-scored without retraining, and the original failure record was retained rather than
overwritten.

\subsection{Inference-policy Audits}
Neither audit changed the released policy, and in both cases the locked confirmation folds were
never opened. Audit~B is the sharper negative result and is summarized in Table~\ref{tab:soup}: on
the 811 calibration cases, weight-averaging the final checkpoint toward the best-validation
checkpoint moved \emph{every} award component in the wrong direction, monotonically in the amount
mixed in. The only favourable movements were the hard-case 5th-percentile Dice ($0.667$ vs $0.664$
for both soups) and, for the 0.50/0.50 soup, the missed-region rate --- neither an award component.

\begin{table}[t]
\centering
\caption{Audit~B (checkpoint soups), calibration folds 0--2, $n=811$; official evaluator, 0 errors,
all six denominators $=811$. \textbf{Inference-policy-selection evidence, not a performance
estimate.} $\theta_f$ = final checkpoint, $\theta_b$ = best-validation checkpoint.}
\label{tab:soup}
\footnotesize
\setlength{\tabcolsep}{4pt}
\begin{tabular}{lcccccccc}
\toprule
& \multicolumn{2}{c}{ET} & \multicolumn{2}{c}{TC} & \multicolumn{2}{c}{WT} & rank \\
\cmidrule(lr){2-3}\cmidrule(lr){4-5}\cmidrule(lr){6-7}
Policy & DSC & HD95 & DSC & HD95 & DSC & HD95 & gain \\
\midrule
\textbf{C0} (released, $\theta_f$) & \textbf{0.8655} & \textbf{13.26} & \textbf{0.9174} & \textbf{5.73} & \textbf{0.9282} & \textbf{5.86} & --- \\
$0.75\theta_f + 0.25\theta_b$ & 0.8641 & 13.72 & 0.9164 & 6.26 & 0.9266 & 6.49 & $-1.000$ \\
$0.50\theta_f + 0.50\theta_b$ & 0.8646 & 13.42 & 0.9171 & 5.96 & 0.9271 & 5.94 & $-1.000$ \\
\bottomrule
\end{tabular}
\end{table}

\subsection{What Has Not Been Measured}
\label{sec:notmeasured}
We state these explicitly rather than leaving them to inference:
\begin{itemize}
  \item \textbf{Final five-model ensemble accuracy.} Every training case was used to train four of
        the five constituent models, so the ensemble cannot be scored without optimistic bias on any
        training case. The OOF numbers in Table~\ref{tab:oof} characterize \emph{single fold models on
        unseen data}, which is the appropriate unbiased proxy, but they are \emph{not} the ensemble's
        score.
  \item \textbf{Official validation-leaderboard performance.} GAT-26 has \textbf{no official
        validation submission and therefore no official rank}: \pending{official validation score and
        rank, if the owner elects to submit}.
  \item \textbf{Genuine-A10G runtime and memory.} All timing and memory figures we hold were measured
        on a different GPU and are explicitly \emph{not} A10G parity:
        \pending{genuine-A10G peak reserved VRAM and seconds per case}.
  \item \textbf{Final container image identity.} \pending{final image tag and digest}.
  \item \textbf{Hidden-test performance.}\\ \pending{hidden-test results, camera-ready}.
\end{itemize}

\section{Discussion}

\subsection{An Explicit Hierarchy of Evidence}
Challenge papers frequently blur distinct quantities. We separate them:
\textbf{(i)} five-fold OOF performance (Table~\ref{tab:oof}) — unbiased for a single fold model on
data it did not see, and the strongest evidence we possess; \textbf{(ii)} inference-policy audit
results (Section~\ref{sec:audits}) — \emph{post-selection} comparisons on development folds, valid
for choosing between policies but not as performance estimates; \textbf{(iii)} the deployed five-model
ensemble — expected to be at least as good as its constituents, but \emph{unmeasurable} on training
data; \textbf{(iv)} official validation performance — not obtained; \textbf{(v)} hidden-test
performance — the only quantity that determines the ranking, and unavailable until the organizers
score the container. Numbers from tier (i) must never be read as tier (iv) or (v).

\subsection{Official Ranking Metrics}
\label{sec:metric}
The organizers state that the Task-3 leaderboard ranking aggregates the \textbf{Dice Similarity
Coefficient} and the \textbf{Normalized Surface Distance (NSD)}, with ranking by summed metric ranks
and pairwise permutation testing to group teams into statistically indistinguishable tiers. Our
internal, pre-registered selection utility used DSC together with the 95th-percentile Hausdorff
distance (HD95), a boundary-error metric that we fixed before training and did not change afterwards.
The Dice half of our utility coincides with the official metric; the HD95 half is an internal
boundary-quality proxy rather than an official ranking metric, and NSD was treated as a diagnostic.
To avoid leaving this as an untested caveat, we re-ran the \emph{entire} pre-existing candidate set
under the official DSC/NSD utility, using the organizers' own evaluation package and configuration,
with the comparison protocol registered in advance of computing any NSD value. The conclusion is
unchanged: no candidate robustly outperforms the released baseline, which is therefore retained and
now validated against the metric the challenge actually ranks on. The correction is not cosmetic ---
enhancing-tumor component cleanup, which the earlier proxy utility rejected because it degraded the
boundary-distance term, improves \emph{both} enhancing-tumor components under the official metric ---
but its advantage remains within sampling uncertainty, so the pre-registered rule retains the
baseline. A fold-0 architecture diagnostic under the official metric likewise moved slightly toward
the larger preset without reaching robustness, and the architecture was not changed.
\ownerinput{confirm at camera-ready the exact NSD tolerance/configuration and final aggregation used
by the organizers --- the released configuration does not pin a tolerance}

\subsection{Limitations}
Group-level splitting mitigates but cannot provably eliminate same-subject leakage. Cohort-stratified
generalization — arguably the core question of a generalizability task — could not be analyzed
because per-case cohort labels were not derivable from the released metadata, so we cannot report
whether performance is uniform across tumor types. Our tail statistics indicate that a small minority
of cases dominate the boundary-error mean, which is the expected failure mode for an unseen-cohort
task. Finally, the negative results of both policy audits mean the system's inference configuration
is deliberately plain; we did not pursue larger ensembles, cascades, or cohort-specific routing.

\subsection{Ethics and Data Governance}
GAT-26 trains only on GoAT-provided data from a recorded random initialization, with no external,
prior-BraTS, private, or pretrained data or weights, consistent with the subchallenge rules.
Validation and test labels are withheld and were never used for training or tuning. We did not
perform pseudo-labeling, transductive adaptation, or leaderboard-driven tuning. Data used in this
work were obtained through the BraTS Challenge project (Synapse ID syn74274097) and are subject to
the challenge data-usage agreement.

\section{Conclusion}
GAT-26 pairs the official region-based nnU-Net ResEnc backbone with a pre-registered, leakage-safe
protocol in which every tunable inference decision was placed behind a frozen, confirmation-gated
policy. The complete five-fold cross-validation is finished and reconciled over all 1{,}351 training
cases with zero evaluator errors, and both inference-policy audits retained the baseline, so the
released configuration is pre-registered rather than post-hoc. The trained models and the frozen
inference policy exist; the final container image, its genuine-A10G acceptance, official validation
performance, and hidden-test results remain \pending{A10G-2, official validation, hidden test} and
will be reported at camera-ready.

\begin{credits}
\subsubsection{\ackname}
\ownerinput{acknowledgements}
% Organizer-mandated acknowledgement sentence, quoted verbatim from the live Challenge Rules
% (Synapse wiki/639585, re-verified 2026-07-28).
Data used in this publication were obtained as part of the Challenge project through Synapse ID
(syn74274097).
\subsubsection{\discintname} The authors declare no competing interests relevant to this work.
\end{credits}

\bibliographystyle{splncs04}
\bibliography{references}

\end{document}
